\documentclass[10pt,a4paper]{amsart}
\usepackage[margin=19mm]{geometry}
\usepackage{amsmath,amssymb,mathtools}
\usepackage[scr=rsfs,cal=euler]{mathalfa}
\usepackage{xfrac}
\usepackage[unicode,hidelinks,bookmarks=false]{hyperref}
\newcommand{\ie}{i.e.\ }
\newcommand{\cf}{cf.\ }
\newcommand{\resp}{respectively\ }
\newcommand{\Subsection}{\subsection}
\providecommand{\nobreakdash}{\nobreak\hbox{-}}
\setlength{\emergencystretch}{2em}
\multlinegap0pt
\usepackage{bm}
\usepackage{bbm}
\usepackage{dsfont}
\usepackage{xspace}
\usepackage{cancel}
\usepackage{mathtools}
\usepackage{amsthm}
\makeatletter
\@ifundefined{lemma}{\newtheorem{lemma}{Lemma}}{}
\makeatother
\title[Novel Dynamics in Models of Angiogenesis with p-Laplacian di]{Novel Dynamics in Models of Angiogenesis with p-Laplacian diffusion}
\author{Wenbo Zhang, Hossein Asgaribakhtiari, Aishwarya Pawar, Rana D. Parshad}
\date{Source: arXiv:2608.07675v1 (CC BY 4.0)}
\begin{document}
\maketitle

\noindent\emph{Source and licence.} Novel Dynamics in Models of Angiogenesis with p-Laplacian diffusion. Version arXiv:2608.07675v1 (CC BY 4.0), distributed under the Creative Commons Attribution 4.0 International licence, as recorded in the arXiv licence metadata for this exact version and shown on the abstract page https://arxiv.org/abs/2608.07675v1. Full rights notice: licenses/arxiv-2608.07675-CC-BY-4.0.txt.\par\smallskip

\noindent\textit{Editorial scope.} Counted unit: the Lemma whose source label is \texttt{lem:lpb} in the article (source0.tex lines 625--634, source label \texttt{lem:lpb}), delivered together with the article's own complete original proof of it (source0.tex lines 637--739, one \texttt{proof} environment of 103 lines). No mathematics is written, repaired or re-derived: statement and proof are the article's own text line for line. Required context retained uncounted: lem:2.1prime (lines 439--465); lem:pc1 (lines 514--532); lem:pc1aa (lines 534--543); eq:ext3dr (lines 554--565). Every definition, display and label of those lines is retained unchanged. Omitted from this package and not redistributed: 1--438, 466--513, 533--533, 544--553, 566--624, 635--636, 740--1346. Nothing inside a retained excerpt is omitted or altered: every retained body is delivered exactly as the article's lines are written, so no omission receipt of any kind accompanies this package. Citations: the retained excerpts carry no citation command at all, so no bibliography entry of the article is transcribed and no reference is redistributed. Reference adaptation: none was needed and none was made; every reference inside the delivered unit resolves inside it, so no unresolved reference or citation is produced. Printed slips kept literal: no printed word, symbol, hypothesis or number of the counted statement or of its proof was altered. Layout and class: the article's own class is \texttt{siamart251216} with packages including lipsum, amsfonts, graphicx, epstopdf, booktabs, array, algorithmic, amsopn; the wrapper is the lane's pinned \texttt{amsart} editorial wrapper at 10pt on A4, which restates the theorem-like environments the retained text needs and adds the article's own macro definitions that the retained text needs (none), with extra wrapper packages (bm, bbm, dsfont, xspace, cancel, mathtools). The delivered numbering is the unit's own. Assets: no retained span contains a figure environment, a graphics-inclusion command, a PDF-page inclusion, a table or a TikZ picture; no image file is copied into this package, the package declares no asset dependency and no image file is redistributed with it. Licence: version arXiv:2608.07675v1 of this article is distributed under the Creative Commons Attribution 4.0 International licence, as recorded in the arXiv licence metadata for this exact version and shown on the abstract page https://arxiv.org/abs/2608.07675v1. A full rights notice is delivered at licenses/arxiv-2608.07675-CC-BY-4.0.txt. Characters: every retained excerpt is pure ASCII, so the delivered engine, XeLaTeX with its default Latin Modern text font, reports no missing character.
\begin{lemma}\label{lem:2.1prime}
Let $\Omega\subset\mathbb{R}^n$ ($n\ge 1$) be a bounded domain with smooth boundary, $1\le p,q\le\infty$, and $D_1,\mu>0$.

\begin{enumerate}
\item[(i)] If $\dfrac{n}{2}\Big(\dfrac1p-\dfrac1q\Big)<1$, then there exists $C=C(D_1,\mu,n,p,q,\Omega)>0$ such that whenever
$\zeta\in C^{2,1}(\bar\Omega\times(0,T))\cap C^0(\bar\Omega\times[0,T))$ for $T\in(0,\infty]$ solves
\[
\begin{cases}
\zeta_t=D_1\Delta\zeta-\mu\zeta+g, & (x,t)\in\Omega\times(0,T),\\
\partial_\nu\zeta=0, & (x,t)\in\partial\Omega\times(0,T),\\
\zeta(x,0)=\zeta_0(x), & x\in\Omega,
\end{cases}
\]
with $g\in C^0(\bar\Omega\times[0,T))$ and $\zeta_0\in W^{1,\infty}(\Omega)$, it holds that
\[
\|\zeta(\cdot,t)\|_{L^q(\Omega)} \le C\Big(\sup_{s\in(0,t)}\|g(\cdot,s)\|_{L^p(\Omega)}+\|\zeta_0\|_{L^\infty(\Omega)}\Big)
\qquad\text{for each } t\in(0,T).
\]

\item[(ii)] Assume that $\dfrac12+\dfrac{n}{2}\Big(\dfrac1p-\dfrac1q\Big)<1$. Then there exists $\tilde C=\tilde C(D_1,\mu,n,p,q,\Omega)>0$ such that if
$\zeta\in C^{2,1}(\bar\Omega\times(0,T))\cap C^0(\bar\Omega\times[0,T))$ with $T\in(0,\infty]$ is a solution of the above problem, then
\[
\|\nabla\zeta(\cdot,t)\|_{L^q(\Omega)} \le \tilde C\Big(\sup_{s\in(0,t)}\|g(\cdot,s)\|_{L^p(\Omega)}+\|\nabla\zeta_0\|_{L^\infty(\Omega)}\Big)
\qquad\text{for any } t\in(0,T).
\]
\end{enumerate}
\end{lemma}

\begin{lemma}
\label{lem:pc1}
    Consider functions $\phi, \sigma \in C^{1}(\Omega), \Omega \subset \mathbb{R}^{n}, n=1,2, |\Omega| < \infty$, $\epsilon > 0$. Then if $p \geq 2$ we have,
    \begin{equation}
        (|\nabla \phi|^{2} + \epsilon)^{\frac{p-2}{2}}|\nabla \phi|^{2} \geq |\nabla \phi|^{p},
\end{equation}


\begin{equation}
        (|\nabla \phi|^{2} + \epsilon)^{\frac{p-2}{2}}|\nabla \phi| \leq C\left(|\nabla \phi|^{p}+1\right),
\end{equation}

and 

\begin{equation}
        (|\phi|^{p-2}\phi - |\sigma|^{p-2}\sigma)\cdot (\phi - \sigma) \geq C |\phi - \sigma|^{p}.
\end{equation}

\end{lemma}

\begin{lemma}
\label{lem:pc1aa}
    Consider functions $\phi, \sigma \in C^{1}(\Omega), \Omega \subset \mathbb{R}^{n}, n=1,2, |\Omega| < \infty$, $\epsilon > 0$. Then if $1< p <2$ we have,
    \begin{equation}
        (|\nabla \phi|^{2} + \epsilon)^{\frac{p-2}{2}}|\nabla \phi|^{2} \geq \left(|\nabla \phi|^{2}+\epsilon \right)^{\frac{p}{2}} - (\epsilon)^{\frac{p}{2}}.
\end{equation}



\end{lemma}

\begin{equation}\label{eq:ext3dr}
\begin{cases}
\displaystyle
(P_{\epsilon})_t = D\,\nabla \cdot
      \Bigl[\bigl(|\nabla P_{\epsilon}|^{2} + \epsilon \bigr)^{\frac{p-2}{p}}\nabla P_{\epsilon}
            - P_{\epsilon}\,\frac{\delta}{(W_{\epsilon}+\beta)(W_{\epsilon}+\gamma)}\, \nabla W_{\epsilon} \Bigr]
      + rP_{\epsilon}\ \left(1-\frac{P_{\epsilon}}{K}\right),
\\[10pt]
\displaystyle
(W_{\epsilon})_t = D_{1} \Delta W_{\epsilon} + \Bigl(\frac{P_{\epsilon}}{1+\nu W_{\epsilon}}-\mu\Bigr)W_{\epsilon},
\end{cases}
\end{equation}

\begin{lemma}
\label{lem:lpb}
Consider \eqref{eq:ext3dr}, then we have that for $p > \frac{3}{2}$, and all other parameters positive, and for initial data $P_{0}(x) \in C^{\omega}(\Omega), 0 < \omega <1, W_{0}(x) \in C^{2}(\Omega)$, chosen sufficiently small $||P_{0}||_{\infty} < C, ||W_{0}||_{\infty}<C$, for any $m>1$, we have,
\begin{equation}
    \int_{\Omega} (1+P_{\epsilon})^{m} dx \leq 
    C.
    \end{equation}
Here $C$ is a pure constant that depends only on the problem parameters.
   
    \end{lemma}

\begin{proof}

Now we use Lemma's \ref{lem:pc1}-\ref{lem:pc1aa}, and whether $1<p<2$ or $p>2$, we obtain,

\begin{eqnarray}
   && \frac{d}{dt}\left(\int_{\Omega} (P_{\epsilon} +1)^{m}dx\right)+\frac{D}{2}m(m-1) \int_{\Omega} (1+P_{\epsilon})^{m-2}|\nabla P_{\epsilon}|^{p} dx  \nonumber \\ 
&\leq& \frac{\delta}{\beta \gamma}\int_{\Omega} (1+P_{\epsilon})^{m-1} \nabla P_{\epsilon} \nabla W_{\epsilon}dx  + m(m-1)(1+(\epsilon)^{\frac{p}{2}})\int_{\Omega} (1+P_{\epsilon})^{m-2}dx \nonumber \\ 
   && + m(r)\int_{\Omega} (1+P_{\epsilon})^{m}dx -\frac{rm}{K}\int_{\Omega} (1+P_{\epsilon})^{m-1}(P_{\epsilon})^{2} dx. \nonumber  \\
        \end{eqnarray}


We also have 
\begin{equation}
\label{eq:mm1}
   \frac{1}{2} \int_{\Omega} (1+P_{\epsilon})^{m-1}2((P_{\epsilon})^{2}+1) dx \nonumber \\ 
\geq \frac{1}{2} \int_{\Omega} (1+P_{\epsilon})^{m-1}(1+P_{\epsilon})^{2}dx,\nonumber 
        \end{equation}
thus 

\begin{eqnarray}
\label{eq:11e}
   && \frac{d}{dt}\left(\int_{\Omega} (P_{\epsilon} +1)^{m}dx\right)+\frac{D}{2}m(m-1) \int_{\Omega} (1+P_{\epsilon})^{m-2}|\nabla P_{\epsilon}|^{p} dx  
   + \frac{rm}{2K} \int_{\Omega} (1+P_{\epsilon})^{m+1}dx\nonumber \\ 
&\leq& \frac{\delta}{\beta \gamma}\int_{\Omega} (1+P_{\epsilon})^{m-1} \nabla P_{\epsilon} \nabla W_{\epsilon}dx  + m(m-1)(1+(\epsilon)^{\frac{p}{2}})\int_{\Omega} (1+P_{\epsilon})^{m-2}dx \nonumber \\ 
   && + m(r)\int_{\Omega} (1+P_{\epsilon})^{m}dx 
   +\frac{rm}{K} \int_{\Omega} (1+P_{\epsilon})^{m-1}dx.
   \nonumber  \\
        \end{eqnarray}

        We state the following lemma,
        \begin{lemma}
        \label{lem:pe21}
            Consider \eqref{eq:ext3dr}. Consider $2>p>\frac{3}{2}$, then we have that 
\begin{equation}
 \int (1+P_{\epsilon})^{m-1}dx \leq  \left(\int (1+P_{\epsilon})^{4m-2}dx\right)^{\frac{1}{4}} \leq \left( \frac{2m-1}{3}\right)C(m) \int_{\Omega} (1+P_{\epsilon})^{m-2}|\nabla P_{\epsilon}|^{p} dx.
        \end{equation}
\begin{proof}
    The proof follows via the embedding of $W^{1,p}(\Omega) \hookrightarrow \hookrightarrow L^{6}(\Omega)$, for $n=2$, $p>3/2$. Also via, $(\int_{\Omega}fdx)^{4} \leq |\Omega|^{3}\int_{\Omega}f^{4}dx$.
\end{proof}


        \end{lemma}
Note via Young's inequality with $\epsilon$ we have,

\begin{eqnarray}
   && \frac{\delta}{\beta \gamma}\int_{\Omega} (1+P_{\epsilon})^{m-1} \nabla P_{\epsilon} \nabla W_{\epsilon}dx  \nonumber \\ 
&\leq& \frac{Dm(m-1)}{2}\int_{\Omega} (1+P_{\epsilon})^{m-2}|\nabla P_{\epsilon}|^{p} dx + C_{3}\int_{\Omega} (1+P_{\epsilon})^{\left(m+\frac{p}{p-1}-2\right)}|\nabla W_{\epsilon}|^{\frac{p}{p-1}} dx. \nonumber \\
        \end{eqnarray}

Here, $\epsilon^{p}=\frac{\beta \gamma}{\delta}\frac{Dm(m-1)}{2}$, which dictates that $C_{3}=\left(\frac{2 \delta}{\beta \gamma Dm(m-1)}\right)^{\frac{1}{p-1}}$.



Inserting the above estimate into \eqref{eq:11e}, and via Holder-Young's inequality, and lemma \ref{lem:pe21} above, we have,


\begin{eqnarray}
   && \frac{d}{dt}\left(\int_{\Omega} (P_{\epsilon} +1)^{m}dx\right)+\frac{Dm(m-1)}{4} \int_{\Omega} (1+P_{\epsilon})^{m-2}|\nabla P_{\epsilon}|^{p} dx  \nonumber \\ 
&\leq&  \left(\frac{2 \delta\theta_{1}}{\beta \gamma D}\right)^{\frac{1}{p-1}}\int_{\Omega} |\nabla W_{\epsilon}|^{\frac{p}{p-1}\theta_{1}} dx + C_{6}\int_{\Omega} (1+P_{\epsilon})^{\left(m+\frac{p}{p-1}-2\right)\left(\frac{\theta_{1}}{\theta_{1}-1}\right)} + m(r)\int_{\Omega} (1+P_{\epsilon})^{m}dx \nonumber \\
&& -\frac{rm}{K}\int_{\Omega} (1+P_{\epsilon})^{m-1}(P_{\epsilon})^{2} dx.  \nonumber \\
        \end{eqnarray}

Here, $\theta_{1}>1$. Note, via Neumann semi-group,

%and is ensured s.t. $\left(m+\frac{p}{p-1}-2\right)\left(\frac{\theta_{1}}{\theta_{1}-1}\right)<m+1$, 



\begin{eqnarray}
\label{eq:wwp}
   && \int_{\Omega} |\nabla W_{\epsilon}|^{\frac{p\theta_{1}}{p-1}} dx  \nonumber \\ 
&\leq&  C_{5}||\nabla W_{\epsilon}||^{\frac{p\theta_{1}}{p-1}}_{\frac{p\theta_{1}}{p-1}} \leq C_{6}||(-\Delta + 1)W_{\epsilon}||^{\frac{p\theta_{1}}{p-1}(a_{2})}_{L^{m}(\Omega)}||W_{\epsilon}||^{\frac{p\theta_{1}}{p-1}(1-a_{2})}_{L^{q}(\Omega)} \leq C_{7}||(-\Delta + 1)W_{\epsilon}||^{\frac{p\theta_{1}}{p-1}(a_{2})}_{L^{m}(\Omega)}.\nonumber \\
 \end{eqnarray}

Next via lemma \ref{lem:2.1prime}, we have,

\begin{equation}
    \int^{t}_{t-1}||(-\Delta + 1)W_{\epsilon}||^{\frac{p\theta_{1}}{p-1}(a_{2})}_{L^{m}(\Omega)}ds \leq \sup_{\tau \in (0,t)}\left(||W_{\epsilon}||^{m}_{L^{m}(\Omega)}\right)^{\frac{p\theta_{1}}{p-1}},
\end{equation}
the mean value theorem for integrals gives the existence of a time $t^{*}$ s.t,

\begin{equation}
    (t-(t-1)||(-\Delta + 1)W_{\epsilon}(t^{*})||^{\frac{p\theta_{1}}{p-1}(a_{2})}_{L^{m}(\Omega)}ds \leq \sup_{\tau \in (0,t)}\left(||P_{\epsilon}||^{m}_{L^{m}(\Omega)}\right)^{\frac{p\theta_{1}}{p-1}},
\end{equation}

since t is arbitrary in $(0,T^{\epsilon}_{max})$, we have that for any $t$ in this time interval,


\begin{eqnarray}
   && \frac{d}{dt}\left(\int_{\Omega} (P_{\epsilon} +1)^{m}dx\right)+ \frac{Crm}{K}\int_{\Omega} (1+P_{\epsilon})^{m+1} dx \nonumber \\ 
&\leq&  \left(\frac{2 (m^{p-1})\delta\theta_{1}}{\beta \gamma D}\right)^{\frac{1}{p-1}}\left(\int_{\Omega} (1+P_{\epsilon})^{m} dx \right)^{\frac{p\theta_{1}}{p-1}} + C_{6}\int_{\Omega} (1+P_{\epsilon})^{\left(m+\frac{p}{p-1}-2\right)\left(\frac{\theta_{1}}{\theta_{1}-1}\right)} \nonumber \\
&& + m(r)\int_{\Omega} (1+P_{\epsilon})^{m}dx . \nonumber \\
        \end{eqnarray}

This enables us to compare to the ODE, for large $m$,
$\frac{dX}{dt} =- X + C_{1}X^{\frac{p\theta_{1}}{p-1}} $.

Choosing $\theta_{1}=2, p=\frac{3}{2}$, we obtain $C_{1} = (\frac{4 \delta}{\beta \gamma D})^{2}$,
Now, uniform boundedness of $X$, is ensured as long as the initial data satisfies, $X(0)=\int_{\Omega} (1+P_{\epsilon}(0))^{m} dx \leq \frac{1}{C_{1}}$. Taking supremum yields,
\begin{equation}
\boxed{||P_{\epsilon}(0)||_{\infty} \leq \left(\frac{\beta \gamma D}{4 \delta}\right)^{2}.}
\end{equation}
\end{proof}

\end{document}
